\section{Evaluation setup}
\label{sec:setup}
\paragraph{Language-model agents.} DeepSeek V4.1 Flash (official API, thinking disabled), MiMo V2.6 Flash (Xiaomi endpoint) and GLM-5.3 Flash (Z.AI endpoint; reasoning mandatory there and capped at 1{,}024 tokens, a recorded asymmetry), each through its maker's endpoint, temperature 0.7, at most 2{,}048 output tokens per call, in the harness of Section~\ref{sec:interface}. Each model first gives a pre-experiment recommendation from the public task only. Its campaigns are branched as in Section~\ref{sec:steps}.

\paragraph{Main baseline: prior-informed local BO.} The numerical baseline uses the same interface, budget and tools. It starts from the fixed reference in one compartment and orthogonal-design perturbations around it in another, adds two staggered compartments at mid-season (under Full feedback chosen with the public process predictor), and plants a second wave when the first crops finish. Its Gaussian process has a season-interaction kernel whose hyperparameters were fitted once on development data and frozen; expected improvement is maximised only inside a box trust region around the best tried policy \citep[TuRBO rules,][]{eriksson2019turbo}, with an initial radius of \res{radius.selected}\% of each range (the largest of three candidates) chosen on development sites by a rule fixed in advance; only tried policies are recommended. This baseline replaced a plain BO after the pilot showed the plain BO scoring below the fixed reference (Appendix~\ref{app:baseline}); the decision referred only to the fixed reference, never to language-model results.

\paragraph{Knowledge asymmetry.} The comparison is not knowledge-symmetric: the baseline starts from the fixed reference and uses a kernel fitted on development data; the models are given neither, and start from their own pre-experiment recommendations. E1 therefore compares general-purpose agents with a domain-informed optimiser.

\paragraph{Research questions.} Three preregistered questions, each answered by one experiment:
\textbf{RQ1} (E1) How good is the final policy a model finds, relative to the baseline with the same interface and budget?
\textbf{RQ2} (E2) Does reading process measurements during the season (Full) improve on seeing only end-of-crop totals (Endpoint)?
\textbf{RQ3} (E3) Are poor recommendations caused by the evidence gathered or by how it is read? A \emph{Reader} here is step 6 of Figure~\ref{fig:framework} taken on its own: a procedure that turns a campaign's final evidence into one recommended policy.

\paragraph{Design and statistics.} \res{sites} unseen test sites; each language model ran two repeats per site, each a branched Full/Endpoint campaign, and the baseline ran two seeds under each condition. The site is the unit: repeats are averaged within a site, and the sample size came from a power analysis on eight pilot sites. The primary family has seven tests with Holm correction: E1, each model's Full final policy minus the baseline's, and E2, Full minus Endpoint for each of the four methods. Each is a site-level paired mean difference, tested by a two-sided sign-flip permutation test ($10^5$ random sign flips) and given a bias-corrected and accelerated (BCa) bootstrap 95\% interval ($10^4$ site resamples); differences below the smallest effect of interest, 2~\eur{} fixed in advance, count as practically negligible. Secondary results, with Benjamini--Hochberg correction over one joint family, are each model's gain over its pre-experiment recommendation, each method against the fixed reference, the between-method differences of feedback gains and E3. Format errors, budget exhaustion and fallback recommendations are the method's own result.

\paragraph{Preregistration and deviations.} Sample size, tests, models, prompts, code and configuration were locked by hash before the run, and the analysis code before any score existed (Appendix~\ref{app:protocol}). Campaigns ran in site order rather than random order; all methods of a site ran in the same window. After the first part of the run, a field-order error in the inputs of the process predictor and the baseline was found; under a new lock, the affected records (all baseline campaigns and \res{repair.branches} language-model Full branches) were rerun from the point of the error, with the baseline radius re-selected by the preregistered rule. Results use the corrected data; the original results are in Appendix~\ref{app:original}.
